\documentclass[11pt]{article}
\usepackage[hmargin=0.9in,vmargin=0.85in]{geometry}
\usepackage{amsmath,amssymb,amsthm}
\usepackage[T1]{fontenc}
\usepackage{lmodern}
\usepackage[hidelinks]{hyperref}
\setlength{\emergencystretch}{1em}
\newtheorem{theorem}{Theorem}
\newtheorem{lemma}{Lemma}
\newcommand{\N}{\mathbb N}
\newcommand{\F}{\mathbb F}
\newcommand{\G}[2]{G(#1,#2;q)}
\title{The Printed Gaussian Coefficient-Sum Sequence\newline
  Is Not Eventually Periodic, Even Modulo Three}
\author{idealistichacker\\\small AI-assisted disproof of conjecture 00000007686}
\date{October 26, 2026}
\begin{document}
\maketitle
\begin{abstract}
Both source languages define $p_q(n)$ as the coefficient sum of the Gaussian
polynomial ${2n\brack n}_q$. That sum is $c(n)=\binom{2n}{n}$, not the ordinary
partition number. The natural-valued sequence is strictly increasing.
More strongly, $c(n)\bmod3$ is not eventually periodic: it has arbitrarily
long zero intervals immediately before powers of three, where its value
is two. This refutes the source's necessary eventual-periodicity clause
at $k=1$, including the plausible implicit coefficient modulus
$p=3,a=1$. The formalization uses actual polynomials, a generated semiring
congruence, an actual polynomial-ring quotient over $\F_3$, and universal
Lucas/zero-interval proofs. No claim about ordinary partition periodicity
or a particular unpublished period formula is made.
\end{abstract}

\section{Exact Scope and the Necessary Clause}
The complete frozen bilingual source is retained separately. Its identity is:
\begin{quote}\small\ttfamily
Commit: b355e42906e6eaf52bc7e9d005e0ef02d8aa7400\\
Blob: 337d3e5e4d97e64cec5d370fa8a3b157ebd77c6f.
\end{quote}
Both languages first mention ordinary partition numbers $p(n)$ modulo $p^a$,
but then \emph{explicitly define} $p_q(n)$ as the coefficient sum of
${2n\brack n}_q$. They assert eventual periodicity in the displayed quotient
$\N[q]/(q-1)^k$, an explicit minimal-period formula, and a further
classical-period comparison. For a sequence $s$ in any fixed quotient, the
necessary clause being refuted is precisely
\[
 \operatorname{EP}(s):\quad
 \exists T\in\N,\ T>0,\ \exists N\in\N,\
 \forall n\ge N,\quad s(n+T)=s(n).
\]
A homomorphism carries this property to its image with the same $N,T$.
It is enough to disprove one such necessary image clause at $k=1$;
neither of the additional period-formula assertions is independently evaluated.
The parameter witness is $p=3,a=1,k=1$, with a prime and positive exponents.
This is not a claim of nonperiodicity for every coefficient modulus:
$c(n)\bmod2$ is zero for $n>0$.

\section{The Actual Gaussian Polynomial and Its Coefficient Sum}
Use the standard nonnegative Gaussian polynomial $G(N,K;q)\in\N[q]$,
characterized by $G(0,0;q)=1$, $G(0,K;q)=0$ for $K>0$, and
\begin{align*}
 G(N+1,0;q)&=1,\\
 G(N+1,K+1;q)&=G(N,K;q)+q^{K+1}G(N,K+1;q).
\end{align*}
The recursion uniquely determines all entries. It gives $G(N,K;q)=0$ for
$K>N$ and $G(N,N;q)=1$. This is an ordinary q-Pascal characterization,
not a sequence declared equal to the desired scalar.
For comparison with the product convention, in $\mathbb Q(q)$ set
\[
 H(N,K)=\prod_{j=1}^{K}\frac{1-q^{N-K+j}}{1-q^j}
 \quad(0\le K\le N).
\]
For $0<K<N$, the ratios to $H(N,K)$ of $H(N-1,K-1)$ and
$H(N-1,K)$ are $(1-q^K)/(1-q^N)$ and $(1-q^{N-K})/(1-q^N)$.
The first plus $q^K$ times the second equals one. Together with the boundary
values, this verifies the same recursion and hence the standard identification.
The Lean definition uses the recursive polynomial characterization; a
separate rational-function product theorem is not advertised as formalized.

For any polynomial $f$, evaluation at one is its finite coefficient sum:
\[
 \sum_j [q^j]f=f(1).
\]
Evaluating q-Pascal at one gives ordinary Pascal recursion and the same
boundary values. Induction in $N$, simultaneously for all $K$, proves
\[
 G(N,K;1)=\binom NK,\qquad
 c(n):=\sum_j[q^j]G(2n,n;q)=\binom{2n}{n}.
\]
There is no evaluation of a $0/0$ product at $q=1$. Both the actual
polynomial recursion and the finite coefficient-sum/evaluation bridge are
proved in Lean. The strict monotonicity of actual central binomial
coefficients also proves $c(n+T)>c(n)$ for every $T>0$, excluding
natural-valued eventual periodicity. This monotonicity is \emph{not} used
to infer modular nonperiodicity.

\section{The q=1 Quotients Are Not Substituted Models}
Natural-coefficient polynomials form a semiring, not a ring. Thus the printed
$\N[q]/(q-1)$ is not silently fed to a commutative-ring ideal API.
The standard semiring presentation is the least additive and multiplicative
congruence generated by $q\sim1$. We prove it equals the kernel congruence
of the actual evaluation homomorphism $E_1:\N[q]\to\N$.

Indeed, $E_1(q)=E_1(1)$, so every generated relation preserves evaluation.
Conversely $q\sim1$ implies $q^j\sim1$. Expanding any polynomial as its
finite sum of monomials gives $f\sim C(f(1))$. Equal evaluations therefore
give congruent polynomials. Constants retain their natural values, so the
source-defined scalar is represented by $C(c(n))$ and cannot acquire a
period in this quotient.

To cover the background's plausible omitted coefficient modulus, use the
genuine ring quotient
\[
 R=\F_3[q]/\langle q-1\rangle.
\]
Its standard evaluation isomorphism $R\simeq\F_3$ sends $C(c(n))$ to
$c(n)\bmod3$. The Lean theorem proves equality of the two actual quotient
classes \emph{if and only if} the corresponding natural remainders agree.
Reduction of the actual polynomial $G(2n,n;q)$ also commutes with evaluation
and has the same class as $C(c(n))$. Thus retaining the Gaussian polynomial
instead of its source-defined coefficient sum does not rescue periodicity.
Reading the printed natural quotient through the ring completion
$\mathbb Z[q]/\langle q-1\rangle$ has the same necessary reduction to
$\F_3$; no separate Lean ring-completion theorem is claimed or needed.

These are the conventional q=1 meanings, not an assertion that the literal
natural-semiring ideal notation is well-formed. Replacing algebraic
subtraction by truncated natural subtraction would impose q=0, a different
repair. Replacing $c(n)$ by ordinary $p(n)$ would change the explicit
definition in both languages. Neither repair is performed here.

\section{Universal Nonperiodicity Modulo Three}
Write $d(n)=c(n)\bmod3$. One-digit Lucas congruences give, for every $m\ge0$,
\begin{equation}\label{digits}
 d(3m)=d(m),\qquad d(3m+1)=2d(m)\pmod3,\qquad d(3m+2)=0.
\end{equation}
For clarity, the last formula does not assume carry-free doubling:
$2(3m+2)=3(2m+1)+1$, whereas $3m+2$ has low ternary digit two;
the low-digit binomial factor is $\binom12=0$.
The first two formulas similarly have low factors $\binom00=1$ and
$\binom21=2$. The formal proof applies the actual Mathlib Lucas theorem.

\begin{lemma}\label{blocks}
For every $m\ge0$, $d(3^m)=2$ and
\[
 d(n)=0\quad\hbox{whenever}\quad 2\cdot3^m\le n<3^{m+1}.
\]
\end{lemma}
\begin{proof}
The power statement follows from $d(1)=2$ and the first formula in
\eqref{digits}. The zero-interval base case is $n=2$.
For the induction step, divide an index in the next interval by three.
Its quotient belongs to the preceding zero interval; its remainder is
zero, one or two. The three formulas in \eqref{digits} then give zero
in every case. This proves the whole interval, not a bounded sample.
\end{proof}

\begin{theorem}
The sequence $\binom{2n}{n}\bmod3$ is not eventually periodic.
\end{theorem}
\begin{proof}
Let any proposed period $T>0$ and starting point $N$ be given. Choose
$m$ with $A=3^m\ge\max(T,N)$ and put $n=3A-T$.
Then $2A\le n<3A$ and $n\ge N$. Lemma~\ref{blocks} yields
\[
 d(n)=0,\qquad d(n+T)=d(3^{m+1})=2,
\]
contradicting the proposed period. In Lean the choice is explicitly
$m=\max(T,N)$, using $m<3^m$; there is no existence assumption hiding
the obstruction. All $T,N$ are quantified.
\end{proof}

Composing this theorem with the actual coefficient-sum and quotient bridges
negates the shared necessary clause at $p=3,a=1,k=1$. Recursive digit
structure is compatible with this nonperiodicity and is not an eventual
fixed additive period. We do not exchange the quantifiers of
coefficient-by-coefficient periods of an unquotiented infinite family.

\section{Formal Scope, Reproduction and Disclosure}
The four mathematical modules are Gaussian, QOnePresentation, ModThree and
PeriodicityNegation. They export 33 public theorems, with all compiled types
and individual axiom footprints audited. The final theorem
\texttt{source\_k\_one\_not\_periodic} gives the natural-presentation and
actual modulo-three quotient negations. A further theorem covers the
coefficient-reduced Gaussian polynomial itself. No forbidden proof escapes
or additional mathematical axioms are used.

Lean is pinned to 4.33.1 and Mathlib to
\begin{quote}\small\ttfamily
0df444a360eaa60ab8c11dca51a86af692955474.
\end{quote}
Only \texttt{propext}, \texttt{Classical.choice} and \texttt{Quot.sound}
occur. Finite exact-integer diagnostics are regression evidence, not the
universal proof. Installed Lean and fixed dependency build caches remain
trust inputs; no full-dependency clean build, empty-machine bootstrap,
cross-platform run or hosted-CI result is implied. Independent AI review
is not human expert endorsement or organizer acceptance. Publication,
priority, scores and awards are not claimed by this mathematical report.
\end{document}
